% !TEX root = ../licenta.tex
\chapter{Introduction}
\label{ch:introduction}

\section{Macro-Environmental Context and Theoretical Motivation}
Traditional macroeconometric forecasting relies heavily on foundational axioms regarding the behavior of time-series data, specifically the assumptions of statistical stability and long-term predictability. Under these strict mathematical assumptions, it is posited that the statistical properties of a given economic indicator—such as its mean, variance, and autocorrelation structure—remain relatively constant over time. However, the global economy is not a closed, sterile laboratory system; it is a highly dynamic, interconnected network constantly adjusting to exogenous shocks, monetary policy shifts, geopolitical conflicts, and rapid technological advancements \cite{ClementsHendry1998}. 

The friction between these rigid mathematical prerequisites and the chaotic reality of global markets is profoundly highlighted during severe macroeconomic crises, such as the 2008 Global Financial Crisis, the European sovereign debt crisis, or the global supply chain collapse induced by the COVID-19 pandemic of 2020. During these periods of unprecedented volatility, historical time-series data suddenly loses its predictive forecasting power. Models trained on decades of stable, pre-crisis data will almost invariably fail to predict post-crisis trajectories because the underlying structural parameters of the economy have fundamentally shifted. This phenomenon is closely related to the famous Lucas Critique \cite{Lucas1976}, which argues that it is mathematically naive to predict the effects of a change in economic policy entirely on the basis of relationships observed in historical data, especially when that data is aggregated across different, incompatible economic regimes.

Given these realities, modern econometricians face a daunting challenge: how to accurately model future economic states when the past is no longer a reliable indicator of the future. The solution necessitates a paradigm shift from static, overarching global models to dynamic, adaptive systems capable of identifying precisely when a regime change has occurred, and subsequently recalculating forecasts based only on the newly established economic reality \cite{Pesaran2004}.

\section{The Universal Benefits of Adaptive Forecasting}
While the primary focus of this thesis is rooted in macroeconomics and econometrics, the mathematical forecasting models and architectural pipelines developed herein have profound benefits that extend far beyond financial markets. Time-series forecasting is a universal mathematical language used to predict the future based on sequential historical data across a multitude of disciplines \cite{Makridakis1998}.

In the field of \textbf{Economics and Finance}, accurate forecasting is the bedrock of fiscal policy, risk management, and algorithmic trading. Central banks rely on inflation and Gross Domestic Product (GDP) forecasts to set interest rates. However, traditional models often fail to account for "Black Swan" events. An adaptive model that can instantly recognize a structural break allows financial institutions to stop incurring significant capital losses by discarding pre-shock data and forecasting based solely on the new market volatility.

In \textbf{Supply Chain Management and Logistics}, forecasting dictates inventory levels, shipping routes, and manufacturing schedules. A sudden disruption, such as a blocked international trade route or a pandemic-induced spike in demand for specific goods, introduces a massive structural break in retail data. An adaptive segmented forecasting system allows logistics companies to immediately identify the break and forecast future demand without the historical data dragging down the new, urgent supply requirements.

In \textbf{Epidemiology and Public Health}, forecasting models are extensively used to predict the spread of infectious diseases. A structural break in epidemiology often corresponds to a viral mutation or the sudden implementation of strict lockdown measures. By detecting the exact point of intervention (the change point), health officials can generate a segmented forecast that proves mathematically whether the intervention successfully altered the exponential trajectory of the infection rate.

In \textbf{Energy Grid Management}, predicting electricity consumption is vital for preventing rolling blackouts and optimizing renewable energy storage. Structural breaks occur due to extreme weather events or sudden shifts in industrial consumption patterns. The ability to detect these shifts and apply localized Moving Averages or Holt-Winters models ensures that the grid remains balanced even under unprecedented stress.

\section{Problem Statement and the Architectural Software Gap}
Despite the urgent cross-industry need for dynamic forecasting mechanisms, there remains a significant architectural gap in the current software ecosystem. Through extensive market research and competitive analysis conducted prior to the development of this project, it became evident that there is no unified product on the market that seamlessly integrates raw econometric modeling, autonomous structural break detection, interactive segmented forecasting, and deterministic Large Language Model (LLM) analytics into a single, decoupled web application.

Existing software solutions generally fall into two disparate, highly polarized categories:
\begin{itemize}
    \item \textbf{Highly Technical Data Science Environments:} Tools such as Jupyter Notebooks, RStudio, and raw Python scripts offer significant computational power. They can execute complex algorithms like the Pruned Exact Linear Time (PELT) algorithm for change point detection \cite{Killick2012} or fit advanced Autoregressive Integrated Moving Average (ARIMA) models \cite{BoxJenkins1970}. However, these environments are entirely inaccessible to non-technical end users. They require extensive programming knowledge, strict library environment management, and a deep understanding of statistical theory, effectively trapping valuable insights within the silos of the data science department.
    \item \textbf{Generic Business Intelligence (BI) Dashboards:} Commercial tools like Tableau, Microsoft PowerBI, and Looker excel in data visualization and democratizing access to interactive charts. Yet, they suffer from a severe deficiency in advanced, on-the-fly statistical modeling. These platforms lack the native, programmatic flexibility to perform automated structural break detection or dynamically calculate partitioned forecasting components without relying on clunky, highly engineered external data pipelines that take weeks for database administrators to deploy.
\end{itemize}

The application developed in this thesis directly addresses this precise market gap. By processing user-uploaded CSV data through an Express.js server, efficiently compiling computationally intensive algorithmic matrix calculations natively within the V8 engine, and rendering the results instantly via a React Virtual DOM, the architecture achieves massive performance benchmarks while maintaining a highly responsive, executive-friendly interface.

\section{Change Points and Segmented Forecasting}
To address the issues of economic shocks highlighted by the Lucas Critique, the application utilizes the concept of Change Point Detection (CPD) \cite{BaiPerron1998} and Segmented Forecasting. 

A \textbf{Change Point} (or structural break) is a specific temporal coordinate in a time series where the fundamental statistical properties of the data abruptly change. This could manifest as a shift in the mean (a sudden jump or drop in asset value), a shift in variance (a sudden period of extreme market volatility), or a shift in the overall trend direction. Identifying these points is critical because fitting a single mathematical model across a change point results in extreme forecasting bias and massive error rates.

\textbf{Segmented Forecasting} is the architectural response to the detection of these change points. Instead of treating the entire history of an economic indicator as a single, continuous dataset, the timeline is treated as a series of distinct "regimes" or segments. When a user interacts with the application, they can visually observe the algorithmically detected change points. If an economic shock is detected, the user can trigger a segmented forecast. The application mathematically truncates the array, permanently discarding the obsolete pre-shock data from memory. The forecasting algorithms are then re-initialized and trained exclusively on the data from the new economic regime. This highly localized approach ensures the forecast is not weighed down by the statistical inertia of an obsolete economic regime.

\section{AI Interpretation: Semantic Context Beyond Mathematics}
The final, and perhaps most innovative, component of this architecture is the integration of Artificial Intelligence to bridge the gap between raw mathematics and human contextual reasoning. Mathematical models are inherently syntactic; they understand the numbers, the variance, and the trajectory, but they have absolutely no understanding of \textit{why} the numbers are moving. A mathematical model does not know what a "pandemic," a "recession," or a "supply chain crisis" is. It only sees a structural break in the temporal variance.

This is where the Large Language Model (LLM) becomes indispensable. By passing the mathematical results into an LLM, we grant the system \textbf{Semantic Understanding} \cite{Korinek2025}. The LLM can explain if the context makes sense beyond what the mathematical model alone can deduce. For example, if the PELT algorithm detects a massive structural break in GDP in Q1 2020, the mathematical model simply calculates a high error rate and adjusts the curve. The LLM, however, possesses vast training data regarding world history and macroeconomics. When prompted with the date of the structural break and the magnitude of the mathematical shift, the LLM can explicitly narrate the overarching economic paradigm to the user.

However, generative AI models, built on probabilistic transformer architectures \cite{Vaswani2017}, are notorious for "hallucinations"—generating confident but factually incorrect mathematical statements. To prevent this, the architecture employs \textbf{Deterministic AI Prompt Engineering}. The AI is never allowed to look at the raw data and guess the forecast. Instead, the Node.js backend executes the strict mathematical proofs, calculates the exact improvement in accuracy (the delta in Root Mean Squared Error), and supplies these pre-calculated values into a hidden heuristic prompt. The LLM is constrained to use the injected results as its baseline, ensuring that its semantic explanation of the context is grounded in verified mathematical results.

\section{Structure of the Thesis}
The remainder of this thesis is organized to provide a comprehensive breakdown of both the mathematical theories and the software engineering principles utilized in the project. Chapter 2 rigorously outlines the theoretical foundations, detailing the pure mathematics behind Moving Averages, Exponential Smoothing, ARIMA models, the PELT structural break detection algorithm, and the mathematical justification for Segmented Forecasting. Chapter 3 dissects the software engineering architecture, including the optimized Node.js V8 execution engine. Chapter 4 presents the full application design and implementation, covering the frontend visualization layer, the backend service architecture, and the deterministic AI integration via the Gemini API. Chapter 5 concludes the research and outlines future development directions.
